% \subsection{Unsupervised Anomaly Detection for Time Series}

% Anomaly detection in time series data has been extensively studied in the literature with various approaches such as statistical methods \cite{paffenroth2018robust}, distance-based methods \cite{chaovalitwongse2007time}, density-based methods \cite{breunig2000lof}, isolation-based methods \cite{liu2008isolation} and one-class classification methods \cite{scholkopf1999support}. The emergence of deep learning led to many novel approaches including Deep-SVDD \cite{pmlr-v80-ruff18a}, GAN-based approaches \cite{li2019mad, Zenati2018AdversariallyLA, zhou2019beatgan, geiger2020tadgan}, reconstruction-based approaches \cite{audibert2020usad, paffenroth2018robust, park2018multimodal, su2019robust, xu2018unsupervised, zhou2019beatgan} or transformers-based approaches \cite{qin2022decomposed, tuli2022tranad} where all of them, mainly rely on unsupervised learning due to the scarcity of labeled data.

% Reconstruction-based models utilize autoencoders to train models with only normal data and detect anomalies by reconstructing a point or a sliding window from test data and comparing them to the actual data according to a reconstruction error.  
% %detect anomalies by assessing the reconstruction error on a test set, having been trained solely on normal data for reconstruction purposes. 
% Variational autoencoders have been applied to this approach, assuming that the training data follows a Gaussian distribution. However, these models often struggle with generalization to other datasets that exhibit different data distributions. On the other hand, GAN-based approaches and one-class classification methods are prone to model collapse during training, leading to instability issues \cite{Arjovsky2017TowardsPM}. Other studies leveraged transformers to efficiently extract temporal information for time series anomaly detection but training transformers demand significant computing resources \cite{zhuang2023survey}.

% % For example, several studies leveraged autoencoders for anomaly detection. Zhao et al. \cite{9964447} employed multiple sequential auto-encoders and training loss balancing for detecting anomalies in multivariate KPIs (Key Performance Indicator). \cite{Johari2022AnomalyDA} introduced a teacher-student-based approach using autoencoders for anomaly detection in Network Functions Virtualizations (NFV) systems. Yin et al. \cite{iot_ad} combined a CNN and a LSTM autoencoder for anomaly detection in IoT time series data. Additionally, other works used GAN-based networks for intrusion detection \cite{gan_ids_araujo} or for KPI anomaly detection \cite{tsagen_wang}.

% Compared to the aforementioned approach, in this work contrastive learning is adopted to enhance the anomaly detection performance and the robustness to data contamination.


%\subsection{Contrastive Learning for Time-Series}


Contrastive Learning has emerged as a prominent self-supervised learning technique demonstrating a great potential across various domains, including computer vision and natural language processing. Traditional contrastive learning models \cite{Caron2020UnsupervisedLO, chen2020simple, He2019MomentumCF} construct positive sample pairs i.e. augmented views of the same instance and negative pairs i.e. augmented views of other instances or a dictionary queue to facilitate representation learning with data augmentation techniques. The effectiveness of these approaches relies on key factors such as data augmentation, efficient sampling of negative pairs, and the use of large batch sizes \cite{chen2020simple}. Recent advancement in contrastive learning architectures such as \cite{Chen2020ExploringSS, grill2020bootstrap} have shown that meaningful representations can be learned without the explicit need of negative pairs or large batch sizes. These techniques have demonstrated remarkable performance in computer vision domain's downstream tasks.

In the domain of time series analysis, recent techniques have been developed to learn representations from time-series data. \cite{franceschi2019unsupervised} introduced a triplet loss and a temporal negative sampling strategy to construct pairs for contrastive learning. \cite{tonekaboni2021unsupervised} leveraged the local smoothness of time series to define neighborhoods and incorporates de-biased contrastive learning. \cite{Eldele2021TimeSeriesRL} utilized temporal and contextual contrastive learning on different views of time-series data generated with weak and strong data augmentation. Similarly, TS2Vec \cite{Yue2021TS2VecTU} employed a hierarchical contrastive learning approach, using augmented context views to capture multi-scale information. Information theory-based adaptive data augmentation is proposed by \cite{Luo2023TimeSC} in their InfoTS method for contrastive learning. \cite{woo2022cost} utilized time domain and frequency domain contrastive losses to learn disentangled seasonal and trend representations for time series forecasting. \cite{Zheng2023SimTSRC} employed a siamese network without negative pairs for time series forecasting. These aforementioned methods have demonstrated robust performance across various downstream tasks, notably in forecasting and classification.

CL techniques have also been applied to anomaly detection in time-series data. \cite{qiu2021neural} employed deterministic contrastive learning with learnable transformations to generate diverse views. \cite{wang2023deep} combined negative sample-free contrastive learning with one-class classification to leverage the representation learning capabilities of contrastive learning and the normality assumption of one-class classification for effective anomaly detection. \cite{kim2023contrastive} also incorporated one-class classification scheme to contrastive learning, utilizing negative data augmentations. \cite{li2023contrastive} utilized temporal transformations to generate anomalies from normal windows, thereby enhancing contrastive learning. \cite{jiao2022timeautoad} employed AutoML techniques to automatically configure the anomaly detection pipeline and tune hyperparameters of the contrastive learning loss, enabling discrimination between original samples and generated negative samples. \cite{DARBAN2025110874} proposed a two-stage framework. The first stage leverages time series representation learning by injecting generated anomalies to learn similar representations for temporally close windows. The final stage classifies window time series based on their nearest and furthest neighbors in the feature space, enhancing anomaly detection.

While some studies suggest that negative pairs are not essential for contrastive learning, we believe that negative pairs, constructed through negative data augmentations, bring some learning knowledge that may enhance AD tasks, unlike their impact on other downstream tasks. In contrast to the aforementioned AD approaches, our approach CATS considers temporal similarity using a temporal loss to better discriminate anomalous time series windows.


% Contrastive Learning has emerged as a prominent self-supervised learning technique demonstrating a great potential across various domains, including computer vision and natural language processing. Traditional contrastive learning models \cite{Caron2020UnsupervisedLO, chen2020simple, He2019MomentumCF} construct positive sample pairs i.e. augmented views of the same instance and negative pairs i.e. augmented views of other instances or a dictionary queue to facilitate representation learning with data augmentation techniques. The effectiveness of these approaches relies on key factors such as data augmentation, efficient sampling of negative pairs, and large batch sizes \cite{chen2020simple}. Recent advancement in contrastive learning architectures such as \cite{Chen2020ExploringSS, grill2020bootstrap} have shown that meaningful representations can be learned without the explicit need of negative pairs or large batch sizes, achieving remarkable performance in downstream tasks.

% In the domain of time series analysis, recent techniques have been developed to learn representations from time-series data. For instance, some methods introduced triplet loss and temporal negative sampling \cite{franceschi2019unsupervised}, while others leveraged local smoothness to define neighborhoods \cite{tonekaboni2021unsupervised}. \cite{Eldele2021TimeSeriesRL} utilized temporal and contextual contrastive learning on different views of time-series data generated with weak and strong data augmentation. Similarly, TS2Vec \cite{Yue2021TS2VecTU} employed a hierarchical contrastive learning approach, using augmented context views to capture multi-scale information. 
% These aforementioned methods have demonstrated robust performance across various downstream tasks, notably in forecasting and classification.

% CL techniques have also been applied to anomaly detection in time-series data. Some studies employed deterministic contrastive learning with learnable transformations \cite{qiu2021neural}, while others combined negative sample-free contrastive learning with one-class classification \cite{wang2023deep}. Other approaches \cite{kim2023contrastive, li2023contrastive, DARBAN2025110874} have utilized temporal transformations to generate anomalies from normal windows and have proposed frameworks that enhance anomaly detection through representation learning.


% While some studies suggest that negative pairs are not essential for contrastive learning, we believe that negative pairs, constructed through negative data augmentations, bring some learning knowledge that may enhance AD tasks, unlike their impact on other downstream tasks. In contrast to the aforementioned AD approaches, our approach CATS considers temporal similarity using a temporal loss to better discriminate anomalous time series windows.

